\begin{workedexample}[Worked Example: Junction temperature of a 5~W device]
Junction temperature follows $T_j = T_a + P\,\theta_{ja}$, where the
junction-to-ambient resistance is the series sum
$\theta_{ja} = \theta_{jc} + \theta_{cs} + \theta_{sa}$ of the junction-to-case,
case-to-sink and sink-to-ambient resistances. A device dissipating $5\ \text{W}$
with $\theta_{jc}=2$, $\theta_{cs}=0.5$ and $\theta_{sa}=3\ ^\circ\text{C/W}$ in
$25\ ^\circ\text{C}$ air reaches
\[ T_j = 25 + 5\,(2+0.5+3) = 52.5\ ^\circ\text{C}. \]
Thermal resistance is the direct analogue of electrical resistance --- power is the
``current'' and temperature the ``voltage.'' Because lifetime falls roughly
exponentially with $T_j$, lowering $\theta_{sa}$ with a bigger heatsink buys
reliability directly.
\end{workedexample}

\begin{keytakeaways}
\begin{itemize}[leftmargin=1.4em,itemsep=2pt]
\item $T_j = T_a + P\,\theta_{ja}$ with $\theta_{ja} = \theta_{jc} + \theta_{cs} + \theta_{sa}$ (series thermal resistances, $^\circ$C/W).
\item Thermal resistance is the electrical analogue of resistance: power is the ``current,'' temperature the ``voltage.''
\item Reliability falls roughly exponentially with $T_j$ (Arrhenius); derating keeps $T_j$ well below the rated maximum.
\end{itemize}
\end{keytakeaways}

\begin{exercises}
\item A device dissipates $8\ \text{W}$ with a total $\theta_{ja}$ of $6\ ^\circ\text{C/W}$ in $30\ ^\circ\text{C}$ ambient. Find $T_j$.
\item Which term of $\theta_{ja}$ do you reduce by bolting on a bigger heatsink?
\item Why derate a device below its maximum rated junction temperature?
\end{exercises}

\furtherreading{Cripps~\cite{cripps}; Gonzalez~\cite{gonzalez}.}
